\section{Introduction}\label{sec:intro}

A developer who sees a qubit end in the wrong state asks which lines of the program determine it. For classical
programs, a backward slice answers that question~\cite{weiser1984program}, and programmers use slices when they
read and debug code~\cite{weiser1982programmers, binkley2007source}. Quantum programs need such help more. Their
intermediate states cannot be inspected without disturbing them, and a gate on one qubit can change another, as in
phase kickback, where a controlled gate changes the phase of the qubit that controls it, even when the two qubits
are never entangled, as Section~\ref{sec:example} shows. Quantum developers debug with print
statements, circuit drawings, simulators and manual reading of the code, and report that these methods do not
scale~\cite{zappin2026challenges}. Bug detectors such as LintQ and QChecker find known patterns by static analysis,
which reads a program without running it, % static-definition
but many faults are wrong outputs that no pattern describes~\cite{paltenghi2024lintq, zhao2023qchecker}. A slice
that shows where such a fault can be is the natural aid.

Slicing has been proposed for quantum programs~\cite{yu2025qpdg, kaul2023uniform, murillo2025roadmap}, but no
proposed dependence rule has been checked against what a quantum program computes. For comprehension this matters
more than for any other use of slicing, because a slice that hides code able to change the qubit misleads its
reader, who will not look where the fault is. Such slices also look right. Fig.~\ref{fig:bv} shows a textbook
program whose classical slice reproduces the faulty program's state and leaves out the faulty line. Over the
\numNPairs{} circuit--qubit pairs of our study, reading a controlled gate the classical way, as an assignment that
leaves its control unchanged, gives a wrong slice for \numCpdgUnsoundPct{} of the pairs, and \numFdCpdgZMatchPct{} of
those wrong slices reproduce the outcome histogram of the whole program, so a developer who checks a slice by its
outcomes would accept them.

This paper asks what a quantum slice must keep, and which dependence rule produces it. We define a slice of a set of
qubits by what it must preserve, their joint state for every input, and an operation as a dependence of those qubits
if a change made right after it can change that state. One rule meets both, the \emph{write-all rule}, under which a
gate writes every qubit it touches, its controls included. On qubits its slice is the backward light cone, whose
soundness compilers already rely on~\cite{murali2019triq, qiskit2025}. What has been missing is the case for it as a
slice, that it keeps every dependence, that nothing else proposed for quantum programs adds a dependence it misses,
and what each smaller or different rule loses. We give the case in proofs and measure the losses in a pre-registered
study of every rule a developer or tool builder would otherwise use, on \numNCircuits{} benchmark circuits with an
exact ground truth, \numFamCircuits{} generated circuits and \numRqThreeUnits{} seeded-fault units, with real faults
from Bugs4Q~\cite{zhao2023bugs4q}.

This paper makes the following contributions.
\begin{enumerate}
\item \textbf{Semantics.} A definition of what a slice of a set of qubits must preserve and of when an operation is a
dependence of those qubits, which separates sound slices from slices that keep every dependence.
\item \textbf{The rule and its guarantees.} The slice of the quantum dependence graph (QDG), which the write-all rule
defines and classical bits extend from a measurement to the operations it conditions, is sound for every input, with
the induction mechanized in Isabelle/HOL, keeps every dependence and distributes over criteria. For each of the
\numGateSignalN{} two-qubit gates of Qiskit's standard library, either qubit can change the other, so no gate can
safely write fewer qubits.
\item \textbf{Three kinds of dependence.} Classical data, classical control and quantum data dependence are all a slice
needs. What looks like quantum control dependence is quantum data dependence, and entanglement is neither necessary
nor sufficient for dependence, so entanglement edges can only add code that cannot matter.
\item \textbf{What the alternatives cost.} The classical program dependence graph misses dependences almost only
through its reading of the control, \numCpdgMissControlPct{} of them. The QPDG of Yu and Zhao~\cite{yu2025qpdg} is
sound, but its slices hold \numYzLcExcessRatio{} times as much code that cannot change the criterion as the QDG's,
whose slices come within \numLcExcess{} points of the smallest slice that keeps every dependence. Pruning a slice by
commutation or by simulating one input makes it smaller but drops dependences, and a state-based slice stays correct
for random inputs on only \numHThreeRate{} of the pairs. Slices narrow where a developer looks but do not improve
mutation-based fault ranking or a language model's localization of faults.
\end{enumerate}

The paper does not build a slicer for a particular language. Its definitions are stated on the circuit a program
builds, Section~\ref{sec:lines} states how a slice reaches the program's source lines, and computing slices at scale
on the source developers write is a separate question. Section~\ref{sec:related} reviews related work and
Section~\ref{sec:example} works through the example of Fig.~\ref{fig:bv}. Section~\ref{sec:core} defines the
semantics and the rules, and Sections~\ref{sec:design} and~\ref{sec:results} present the study. The paper ends with
the implications for comprehension tools, the threats to validity and the conclusion. The data, scripts, protocol
and proofs are available in an anonymous replication package.\TODO{anonymous package link}
